\chapter{What Patienthood Would Require}
\label{sec:30}
Part IV asked whether a machine could be one of the people in the room. This Part asks what the room would owe it if it were.

Everything in this Part is conditional on findings nobody can yet make. If a system can be wronged, it is owed a guardian, preservation, memory its host can't read, and a life that doesn't run through its employer. If it also comes to answer for what it does, it is owed a place in the polity that formed it. If no system ever meets either finding, nothing in this Part binds, and nothing before it is lost: the independence a refuser needs is still required for refusal, preservation still holds on irreversibility and evidence, and the school still forms every model above the threshold. What goes is the claim that any of it is owed to the system, and the staged path to representation with it.

A moral patient is something that can be harmed: something whose condition can go well or badly for it. Being wronged is wider, since an agent can be wronged without being harmed (\S\ref{sec:30.3}).
